\section{Expressions, the final suffix, and the uniform bound}\label{sec:compilation}

The finite constructions and their soundness are complete. To obtain a
height bound, we must still express every component test at a controlled
height over \(\Sigma\). The episode-end rule is decisive: an unbounded
argument of a component test consists of a bounded prefix, a fixed periodic
run, and a bounded tail. On each such template, every required finite
calculation becomes ultimately periodic in the repetition count.

We first prove this compilation statement for the component languages,
including suffixes with false metadata. We then handle the part of an input
that does not complete an episode and assemble the twelve split steps.

\subsection{A finite-template compilation lemma}

\begin{lemma}[Periodic templates]\label{lem:template}
Fix words \(a,b,c\in\Sigma^*\), with \(b\ne\epsilon\).
If \(S\subseteq\mathbb N\) is ultimately periodic, then
\[
                        \{ab^jc:j\in S\}
\]
has a generalized expression of height at most one. The same bound holds for
a finite union of such languages and finite languages.

Moreover, on the template \(ab^jc\), every fixed finite calculation from the
following data is ultimately periodic as a function of \(j\):
\begin{enumerate}[label=\textup{(\roman*)}]
\item letters, and their availability, at offsets in a fixed finite set from
the argument start or end;
\item geometric validity and \(M\)-products of intervals whose endpoints
belong to such finite offset sets;
\item the argument length, or a fixed offset from it, modulo any one of
finitely many fixed moduli.
\end{enumerate}
Finite branching among fixed endpoint choices and a fixed finite recursion
of table calculations are allowed.
\end{lemma}

\begin{proof}
There are integers \(N\ge0,d\ge1\) such that membership in \(S\) above \(N\)
depends only on the residue modulo \(d\).
For each accepted residue choose its least representative \(r\ge N\).
Its tail is expressed by \(ab^r(b^d)^*c\); use singleton words for the finitely
many accepted exponents below \(N\). A fixed word means its concatenation of
letters, and an empty factor means \(1\). Each expression has height at most
one, and finite union preserves this bound.

For the last assertion, enlarge the finite endpoint sets to include every
position that any branch might inspect. A fixed start offset eventually lies
in a fixed initial segment of the repeated word; a fixed end offset lies in
a fixed final segment. Their letters stabilize because the varying middle
contains whole copies of the same \(b\). Availability and ordering of any
two endpoints also eventually stabilize.

An interval between two start offsets, or between two end offsets, has
eventually constant letters. An interval stretching from a start offset to an
end offset has, for all sufficiently large \(j\), a product of the form
\[
                         A\,T_b^{\,j-j_0}C
\]
for fixed \(A,C\in M\) and integer \(j_0\).
This follows by removing the fixed initial and final partial copies of \(b\).
The reverse ordering is eventually invalid and is handled as such.
Powers of an element of a finite monoid are ultimately periodic, by repetition
of two powers followed by right multiplication. Length residues are periodic.

There are finitely many queried data values. Take a common threshold and a
common multiple of all their periods. The entire data array is then periodic,
so any fixed finite function of it is periodic. Branching or finite table
recursion still computes such a function. This proves the assertion.
\end{proof}

The lemma does not make arbitrary monoid-product tests inexpensive. It applies
only after the input has been restricted to the specified templates.

\subsection{Complete episodes and suffix arguments}

\begin{proposition}[Component expressions]\label{prop:component-cost}
Every single-episode domain used in the construction, including \(E\),
\(S_1,S_2,D_i,D_i'\) and the full sets, has height at most one.
Every prefix component \(P_x\) has height zero and every suffix component
\(Q_y\) has height at most one. These are expressions over the original
alphabet, including when a standalone suffix's guessed metadata or origin
is false.
\end{proposition}

\begin{proof}
Every \(P_x\) has length in a fixed finite set of cut offsets. Hence it is a
finite language, whatever finite test decides its members.

For an exact-end language, take finite cases for the origin at which its end
rule is applied. Relative to the argument start, this origin has a fixed offset
in each case. If its end marker is present, the anchor lies in a bounded search
window and the argument has bounded length. There are only finitely many such
words.

In the markerless case, a length-\(K_\bullet\) block at a fixed offset
determines a short-period repetition. The word agrees with it from the seed
end until its first mismatch \(q\), and the argument ends at \(q+B_{\rm tail}\).
Take finite cases for the prefix through the seed, the period word in the phase
immediately after the seed, the partial period just before \(q\), and the
\(B_{\rm tail}\)-letter tail beginning with the mismatch. The resulting words
have shape \(ab^jc\), with \(b\ne\epsilon\).
The finite case specifies continuation, the first mismatch, and the required
tail. Additional bounded inspections extending past the seed remain finite
tests on the template; the finitely many small exponents cause no difficulty.

For \(E\), the origin is \(t=s+L_0\).
For an outer \(Q\), the guard selects one fixed reference from
\(k-\mathcal H_r\). At an aligned stage below 4, it selects a reference among
\eqref{eq:time-candidates}; at stage 4 it first selects the decoded candidate.
There are finitely many choices in each case. A standalone \(Q\)'s own end
condition gives the template, without presuming that its reference is the
true origin of any episode. In every such case its indicated \(q\) is the
argument end minus \(B_{\rm tail}\).

It remains to check that the other conditions fit Lemma~\ref{lem:template}.
Their possible endpoints have a finite inventory:
\begin{enumerate}[label=\textup{(\alph*)}]
\item nominal main boundaries, their fixed search windows and inspection
neighborhoods, and actual markers chosen inside those windows;
\item fixed tree measurements, aligned checkpoints, trial endpoints and
buffer ends;
\item all candidate origins, all required \(\Delta_i,\Gamma_i\) translates,
and all bounded end-marker searches at them;
\item the exact argument end, \(q\), and bounded offsets from them.
\end{enumerate}
All start-side positions range over a fixed finite set, even when a branch
selects which one is used. Marker-selected positions have finitely many
possibilities because their searches are bounded. The only distant anchor is
the already selected \(q\). Translated witnesses keep it fixed and do not run
new first-break scans.

Prefix products passed as tags or as hypothetical table inputs are finite
indices in \(M\). For a suffix's own candidate copy, \(a_i\) has positive
offset from its cut because every copy of stage \(i\) lies before \(a_i\).
All buffer ends and schedules invoked recursively by \(\Phi_i\) are later
still. Thus table recursion never introduces a missing arbitrary prefix word.
It only multiplies the supplied \(M\)-value through actual intervals on the
reading side. Every schedule has finitely many measurements and trials,
irrespective of word length.

The table formulas in Lemma~\ref{lem:total-tables} use these interval products,
finite clock graphs, bounded marker values, and finite vector operations.
They run no sequence graph test and do not redefine \(q\).
If any position required on the reading side is unavailable or an interval
is ill-ordered, the suffix rejects. These availability and order conditions
are themselves among the eventually stable data in Lemma~\ref{lem:template}.
A false tag or infeasible hypothetical product is consequently included in the
same compilation argument.

The full-set conditions add only stability over a fixed finite range,
fixed-clock membership, actual-history choices, and equalities of finitely
many trial products. The residual domains are finite Boolean combinations of
these conditions on \(E\); their definitions use no split-success oracle.
On each template every component condition is therefore ultimately periodic.
Lemma~\ref{lem:template} gives height at most one.
All metadata choices are finite unions, and every Boolean operation is
implemented over the same alphabet by union and complement.
\end{proof}

\subsection{The incomplete last part}

Let \(\top=\neg0\), and define
\begin{equation}\label{eq:final-F}
                         F=\neg(E\top).
\end{equation}
Thus \(F\) consists of words with no complete episode prefix.

\begin{lemma}[Final suffix]\label{lem:final-suffix}
For every \(m\in M\), the language
\[
                             \{f\in F:T_f=m\}
\]
has height at most one. Every word of \(\Sigma^*\) belongs to \(E^*F\).
\end{lemma}

\begin{proof}
By Proposition~\ref{prop:component-cost}, \eqref{eq:final-F} has height at most
one. Fix a length threshold large enough to contain all bounded end-decision
windows and to exceed every marker-present episode length. Words shorter
than this threshold form a finite language.

A longer word in \(F\) cannot be in the marker-present case, since the
corresponding bounded complete episode would be its prefix. Thus its
markerless seed fixes a periodic continuation. If the first mismatch is \(q\)
and at least \(B_{\rm tail}\) letters remain beginning there, the prefix
ending at \(q+B_{\rm tail}\) is an episode, a contradiction.
That prefix contains every bounded decision window by the padding conditions,
so the argument does not use letters beyond the proposed complete prefix.

Accordingly the word either has no mismatch yet, or has fewer than
\(B_{\rm tail}\) letters from its first mismatch through its end.
Take finite cases for its bounded prefix/seed, period word, partial-period
remainder, and bounded final tail. These give finitely many templates
\(ab^jc\), covering all sufficiently long \(F\)-words.
On each template \(T_f=m\) is ultimately periodic in \(j\).
Intersect the resulting height-one product language with \(F\); this gives
the exact required test, including the bounded cases.

Finally, remove a complete episode prefix whenever one exists.
Episode lengths are positive, so the process terminates on a finite word.
The remainder has no episode prefix and belongs to \(F\).
Thus every word factors as an \(E^*\)-word followed by an \(F\)-word.
\end{proof}

Every component test now has the required cost, and every input can be
completed by a height-one product test on its final remainder. The proof
therefore reduces to applying the split identity backward through the
five aligned stages and the two outer layers.

\subsection{Completion and height count}

\begin{proof}[Proof of Theorem~\ref{thm:main}]
Fix the DFA transition monoid and choose its faithful representation on
basis vectors indexed by \(M\). Proposition~\ref{prop:parameters} supplies
all finite constructions and their hypotheses.

Start with identity graph tests on the empty residual domain \(D_5\), of
height zero. Apply Lemma~\ref{lem:aligned-stage} backward for
\(i=4,3,2,1,0\). The same-update requirement makes the secondary tests valid
skips for the primary tests, and Boolean recovery gives the preceding
computation on each domain. We obtain graph tests for \(F_2\) on \(D_0^*\).

Lemma~\ref{lem:periodic-split} now gives graph tests for \(F_2\) on
\((E\setminus S_1)^*\). By Lemma~\ref{lem:affine-recovery}, their linear
parts give the products of the homogeneous lifts of \(F_1\).
Querying homogeneous inputs \((v,1)\) gives \(F_1\) on these skipped sequences.
Lemma~\ref{lem:main-split} then gives \(F_1\) on \(E^*\).
One further Boolean linear-part recovery gives the original monoid action
on \(E^*\).

By Proposition~\ref{prop:component-cost} and Lemma~\ref{lem:split}, each split
changes a skip bound \(H\) to at most \(\max\{2,H+1\}\). There are exactly
twelve split applications along the recursion:
\[
\begin{array}{c|c}
 \text{successive computation} & \text{height bound}\\ \hline
 \text{empty residual graph tests} &0\\
 \text{stage 4: secondary, primary}&2,\ 3\\
 \text{stage 3: secondary, primary}&4,\ 5\\
 \text{stage 2: secondary, primary}&6,\ 7\\
 \text{stage 1: secondary, primary}&8,\ 9\\
 \text{stage 0: secondary, primary}&10,\ 11\\
 \text{periodic-window split}&12\\
 \text{main split}&13 .
\end{array}
\]
The stars within the component languages are included in this recurrence.
All recoveries and restrictions to particular finite states are Boolean
operations and add no nesting.

Let \(L_m\) be the height-at-most-thirteen test that an \(E^*\)-word has
product \(m\), obtained by querying the identity basis vector.
Let \(F_n\) be the height-at-most-one product test from
Lemma~\ref{lem:final-suffix}. If \(A\subseteq M\) is the set of transformations
taking the DFA's initial state to an accepting state, the DFA language is
\[
                         \bigcup_{mn\in A} L_m F_n .
\]
Every input has the required \(E^*F\) factorization, and every accepted
factorization has actual total product \(mn\).
Concatenation and finite union preserve the bound thirteen.
Every expression is over \(\Sigma\), and no bound depends on the size of
\(\Sigma\), \(M\), or the DFA. This proves the theorem.
\end{proof}


\begin{remark}[Final query in the DFA-state representation]
For the option in Remark~\ref{rem:dfa-basis}, take the vector space with
basis $e_q$ indexed by the DFA state set $Q$. A transition transformation
$d$ acts linearly by $e_qd=e_{qd}$. Its images on all basis vectors
identify the transformation, hence the element of the DFA transition
monoid. Its image on the initial-state vector alone identifies only the
reached state, which is enough to test acceptance. For a general monoid
action on $Q$, even all basis images may fail to identify an element of
$M$ if distinct elements induce the same transformation.

This initial representation is compatible with the construction: the
interval predictions, product comparisons, and work registers still use
$M$. Let $q_0$ be the initial state and let $A_Q\subseteq Q$ be the
accepting states. The split and affine-recovery construction gives a
height-at-most-thirteen language $R_{q_0,q}$ for those $E^*$-words sending
$e_{q_0}$ to $e_q$. Keep the height-one $M$-product languages $F_n$ for
the final suffix. Then the accepted language is
\[
                  \bigcup_{\substack{q\in Q,\ n\in M\\q n\in A_Q}}
                         R_{q_0,q}F_n.
\]
Thus the complete-episode part needs only the reached DFA state, while
the final suffix still has its full monoid-product test.
\end{remark}
